\pdfoutput=1
\documentclass[11pt]{article}

% -----------------------------------------------------------------------------
% Draft skeleton for a professional paper.
% Purpose: short main text + modular appendices, with bullets/single sentences
% that can be expanded into full paragraphs later.
% -----------------------------------------------------------------------------

\usepackage[margin=1.08in]{geometry}
\usepackage[T1]{fontenc}
\usepackage[utf8]{inputenc}
\usepackage{lmodern}
\usepackage{microtype}
\usepackage{amsmath,amssymb,amsthm,mathtools}
\usepackage{booktabs,longtable,array}
\usepackage{enumitem}
\usepackage{xcolor}
\usepackage{hyperref}
\usepackage[nameinlink,capitalize]{cleveref}
\usepackage{listings}

\hypersetup{
  colorlinks=true,
  linkcolor=blue!60!black,
  citecolor=blue!60!black,
  urlcolor=blue!60!black,
  pdftitle={Local Proof Architectures for Critical PDE Regularity},
  pdfauthor={Guillem Duran Ballester},
  pdfsubject={Proof architecture, critical PDE, Navier--Stokes, LLM-auditable proof design},
  pdfkeywords={proof architecture, state-space stratification, fault containment, critical PDE, Navier--Stokes, LLM audit}
}

% Theorem environments.
\theoremstyle{plain}
\newtheorem{theorem}{Theorem}[section]
\newtheorem{proposition}[theorem]{Proposition}
\newtheorem{lemma}[theorem]{Lemma}
\newtheorem{corollary}[theorem]{Corollary}
\theoremstyle{definition}
\newtheorem{definition}[theorem]{Definition}
\newtheorem{example}[theorem]{Example}
\theoremstyle{remark}
\newtheorem{remark}[theorem]{Remark}

\crefname{theorem}{Theorem}{Theorems}
\Crefname{theorem}{Theorem}{Theorems}
\crefname{proposition}{Proposition}{Propositions}
\Crefname{proposition}{Proposition}{Propositions}
\crefname{lemma}{Lemma}{Lemmas}
\Crefname{lemma}{Lemma}{Lemmas}
\crefname{corollary}{Corollary}{Corollaries}
\Crefname{corollary}{Corollary}{Corollaries}
\crefname{definition}{Definition}{Definitions}
\Crefname{definition}{Definition}{Definitions}
\crefname{example}{Example}{Examples}
\Crefname{example}{Example}{Examples}
\crefname{remark}{Remark}{Remarks}
\Crefname{remark}{Remark}{Remarks}

% Useful macros.
\newcommand{\NS}{Navier--Stokes}
\newcommand{\R}{\mathbb{R}}
\newcommand{\N}{\mathbb{N}}
\newcommand{\cI}{\mathcal I}
\newcommand{\cC}{\mathcal C}
\newcommand{\cG}{\mathcal G}
\newcommand{\cS}{\mathcal S}
\newcommand{\cR}{\mathcal R}
\newcommand{\eps}{\varepsilon}
\newcommand{\dd}{\,d}
\newcommand{\todo}[1]{\textcolor{red!70!black}{\textbf{TODO:} #1}}
\newcommand{\capsule}{\mathfrak C}
\newcommand{\branch}{\mathfrak b}
\newcommand{\ledger}{\mathfrak L}

\setlist[itemize]{leftmargin=1.6em,itemsep=0.25em,topsep=0.35em}
\setlist[enumerate]{leftmargin=1.8em,itemsep=0.25em,topsep=0.35em}
\renewcommand{\arraystretch}{1.16}

\lstdefinestyle{auditstyle}{
  basicstyle=\ttfamily\small,
  breaklines=true,
  frame=single,
  columns=fullflexible,
  backgroundcolor=\color{gray!5}
}
\lstset{style=auditstyle}

\newenvironment{takeaway}
{\par\medskip\noindent\begin{minipage}{\linewidth}\hrule\vspace{0.4em}\textbf{Takeaway.}\ }
{\vspace{0.4em}\hrule\end{minipage}\par\medskip}

\title{\textbf{Local Proof Architectures for Critical PDE Regularity}\\
\large State-Space Stratification, Fault Containment, and LLM-Auditable Proof Design}
\author{Guillem Duran Ballester\\
\small \texttt{[email / affiliation placeholder]}}
\date{Working draft --- \today}

\begin{document}
\maketitle

\begin{abstract}
This paper proposes a verification-aware proof architecture for critical PDE regularity problems based on local state-space stratification.  The framework replaces a monolithic contradiction proof by a finite directed graph of proof capsules, each with explicit hypotheses, outputs, retained obstructions, failure routes, and audit tests.  A candidate singular branch is routed through scale-invariant local tests until it reaches either a closed terminal stratum or a named open obstruction.  By isolating the singularity into compact local strata, the architecture bypasses the need for global critical-norm estimates or whole-space Liouville theorems: nonlocal pressure tails at spatial infinity are walled off, and local Calder\'on--Zygmund control on compact cylinders becomes sufficient for the active branch.  The framework treats a PDE proof as a fault-tolerant software system whose logical structure is an abstract syntax tree suitable for direct LLM processing and eventual formal verification in systems such as Lean~4.  The three-dimensional incompressible \NS{} equations are used as a case study, with Caffarelli--Kohn--Nirenberg concentration, Type I/Type II separation, ancient-profile routing, repaired gauges, scale-cost mechanisms, and terminal profile exits serving as motivating examples.
\end{abstract}

\paragraph{Status of this draft.}
This document is a paper skeleton: every section has the final intended role, but most entries are written as one-sentence bullets or theorem templates to be expanded later.

\paragraph{Core principle.}
The paper studies proof architecture and verification structure by isolating every analytic step into an independently auditable capsule recorded in an explicit ledger.

\tableofcontents

% =============================================================================
\section{Introduction}
\label{sec:introduction}

\subsection{Motivation: why monolithic PDE proofs are fragile}
\begin{itemize}
  \item Critical PDE regularity arguments often fail for architectural reasons: a single misunderstood compactness limit or pressure normalization can silently contaminate dozens of later pages.
  \item Monolithic contradiction proofs force local facts, global estimates, and rigidity arguments into one fragile trunk, so one bad step on page 84 can snap the entire proof.
  \item The real bottleneck is not only human memory; it is a proof format that allows hidden global dependencies to leak across the manuscript without typed interfaces.
  \item LLM-assisted proof search exposes this architectural flaw even more sharply, because a model can verify a local derivation while still missing the global estimate leak that makes the overall argument unsound.
\end{itemize}

\subsection{The Curse of Globality}
\begin{itemize}
  \item Classical regularity programs often try to impose global constraints such as whole-space critical-norm bounds on equations whose dynamics still carry nonlocal terms, most notably pressure.
  \item Those global bounds accumulate interacting cutoff errors, pressure tails, and compactness losses across many scales until the bookkeeping itself becomes the dominant difficulty.
  \item A proof built around global $L^3$ control, global profile decompositions, or whole-space Liouville inputs is therefore structurally predisposed to catastrophic coupling between distant parts of the argument.
  \item The curse of globality is that every new estimate must remain compatible with the entire ambient space at once, even when the singularity mechanism is fundamentally local.
\end{itemize}

\subsection{The Modular/Local Paradigm}
\begin{itemize}
  \item The proposed alternative is to replace global proof trunks by finite state-space dichotomies driven only by local, scale-invariant observables on compact parabolic cylinders.
  \item In this paradigm, global estimates are a relic of monolithic proof structure rather than a mathematical necessity: each branch only needs the local bounds generated inside its own state.
  \item The resulting proof is fault-tolerant and strictly bounded, because every local estimate has a finite area of effect and every branch can be audited without importing the entire ambient space.
  \item This is the same conceptual move that makes modern software reliable: isolate modules, define interfaces, and forbid undeclared dependencies from leaking across the system.
\end{itemize}

\subsection{The proposed alternative}
\begin{itemize}
  \item We organize a regularity proof as a finite routed state space rather than as one uninterrupted global argument.
  \item Each edge in the state-space graph is a proof capsule with a strict local contract: inputs, outputs, retained invariant, failure routes, and audit obligations.
  \item The architecture compiles a continuous infinite-dimensional PDE problem into a discrete mathematical API whose nodes and edges can be checked one local interface at a time.
  \item A failure of a capsule does not erase the whole architecture; it reopens one edge or terminal node and leaves the verified subgraph meaningful.
\end{itemize}

\subsection{Main contributions}
\begin{itemize}
  \item We define proof capsules, retained obstructions, routing graphs, terminal closures, and fault containment.
  \item We elevate elimination of hidden global estimates to a design principle: every branch is expressed through local observables, local gauges, and local closure capsules.
  \item We prove an abstract routed-closure theorem: if every branch is routed and every terminal node is closed against a retained obstruction, then the root branch is impossible.
  \item We propose an LLM-audit protocol for scaling, sign, pressure, compactness, dependency, and failure-route checks, tuned to the local reasoning window of current models.
  \item We introduce a mathematical API that compiles continuous, infinite-dimensional PDE dynamics into discrete local nodes that LLMs and formal systems can navigate without exponential blow-up of the search space.
  \item We give the \NS{} singularity-routing program as a detailed case study and stress test for machine-auditable proof design, with the routing graph functioning as the proof's logical AST.
\end{itemize}

\begin{takeaway}
The novelty is not a single new estimate; the novelty is a proof architecture in which every estimate is local, every failure of an estimate has a semantic destination, and every terminal contradiction is tied to a retained obstruction.
\end{takeaway}


\subsection{Relationship with current literature}
\begin{itemize}
  \item A template for concentration compactness arguments, but with the concentration mechanism promoted from one lemma inside a long proof to the front door of an explicit routing graph.
  \item Similarities to profile decomposition and Littlewood--Paley methods, except that those tools are deployed only on specific strata rather than across the entire state space at once.
  \item Compatibility with partial regularity theory and imported rigidity theorems, which now appear as typed routing endpoints rather than as diffuse background assumptions.
\end{itemize}

% =============================================================================
\section{Overcoming the Global Estimate Bottleneck}
\label{sec:global-bottleneck}

\subsection{Analytic surgery instead of global smoothness enforcement}
\begin{itemize}
  \item The guiding analogy is Perelman's Ricci flow surgery: instead of demanding that the full evolution remain globally well behaved in one stroke, we isolate the dangerous region and classify the singular behavior into finitely many irreducible atoms.
  \item In the \NS{} setting those atoms are compact cores, rough cores, multibubbles, cascades, terminal profiles, and related local descendants.
  \item The proof architecture is therefore surgical rather than monolithic: singular behavior is not denied abstractly, but routed into a finite taxonomy where each atom has its own closure capsule.
\end{itemize}

\subsection{Walling off spatial infinity}
\begin{itemize}
  \item The architecture works on compact parabolic cylinders with explicit cutoff functions, so only the immediate area of effect of the candidate singularity enters the active calculation.
  \item Once the proof is routed into a local stratum, spatial infinity is walled off rather than globally integrated into every later estimate.
  \item This is the structural reason the program can discard whole-space critical-norm control: the branch only pays for the geometry that survives inside its local window.
\end{itemize}

\subsection{Covariant gauge observers and local pressure control}
\begin{itemize}
  \item The pressure term is handled by local gauge fixing and covariant observers adapted to the active branch rather than by a single global reconstruction formula imposed on all scales at once.
  \item Repaired gauges flatten the local connection enough to expose explicit modulation coefficients, while local Calder\'on--Zygmund estimates control the pressure contribution inside the compact cylinder actually used by the branch.
  \item This local observer viewpoint prevents pressure tails at infinity from leaking back into every estimate as a hidden global error source.
\end{itemize}

\subsection{Why this geometry matches machine reasoning}
\begin{itemize}
  \item LLMs and formal tactic engines struggle when asked to reason over a large topological proof trunk whose active hypotheses are spread across many global layers.
  \item They perform far better on compact local tasks such as a cutoff estimate, a local Sobolev embedding, a Calder\'on--Zygmund bound, or a compactness passage inside a fixed cylinder.
  \item The architecture is therefore mathematically stronger and computationally better aligned: it redesigns the proof so that the dominant reasoning objects already fit the cognitive profile of machine auditors.
\end{itemize}

% =============================================================================
\section{Abstract Routed Proof Architectures}
\label{sec:abstract-architecture}

\subsection{Branches, states, and observables}
\begin{definition}[Branch]
A branch is a candidate counterexample trajectory together with the local data retained after all preceding routing steps.
\end{definition}

\begin{definition}[State]
A state is a named local regime characterized by scale-invariant observables, compactness data, and declared failure alternatives.
\end{definition}

\begin{itemize}
  \item Examples of observables include critical concentration, local energy, pressure oscillation, terminal mass, scale drift, profile count, and tail defect.
  \item States are defined purely by local, scale-invariant observables on compact windows; they do not depend on boundary data at spatial infinity.
  \item A state should be stable under the rescalings, gauges, and subsequences used before its closure theorem is invoked.
\end{itemize}

\subsection{Retained obstructions}
\begin{definition}[Retained obstruction]
A retained obstruction is a nonzero scale-invariant quantity attached to a branch and required to survive every allowed extraction, rescaling, gauge change, or descendant operation.
\end{definition}

\begin{itemize}
  \item The retained obstruction is the object contradicted by terminal rigidity or vanishing.
  \item Typical examples are positive CKN concentration, compact local $L^3$ mass, nonzero terminal profile mass, or a nonvanishing localized core.
\end{itemize}

\subsection{Proof capsules}
\begin{definition}[Proof capsule]
A proof capsule is a tuple
\[
  \capsule=(H,O,I,F,A),
\]
where $H$ is the input contract, $O$ is the output contract, $I$ is the retained obstruction, $F$ is the list of named failure routes, and $A$ is the audit checklist.
\end{definition}

\begin{itemize}
  \item A proof capsule is a mathematical API contract: $H$ is the pre-condition, $O$ is the post-condition, and the retained obstruction plus failure routes specify what is preserved if the edge succeeds or how the branch is retyped if it does not.
  \item This strict typing is what prevents hidden global-estimate leaks: every theorem invocation must expose the exact local data it consumes and the exact local data it exports.
\end{itemize}

\begin{center}
\begin{tabular}{p{0.22\linewidth}p{0.68\linewidth}}
\toprule
Capsule field & Meaning \\
\midrule
Input contract $H$ & Exact hypotheses the capsule is allowed to use. \\
Output contract $O$ & The local conclusion produced if the capsule succeeds. \\
Retained obstruction $I$ & The invariant that must survive the capsule. \\
Failure routes $F$ & Named alternatives entered if the capsule hypotheses fail. \\
Audit checklist $A$ & Scaling, sign, compactness, pressure, dependency, and locality tests. \\
\bottomrule
\end{tabular}
\end{center}

\subsection{Routing graphs}
\begin{definition}[Routing graph]
A routing graph is a finite directed graph $G=(V,E)$ with a root state, edge labels given by proof capsules, and terminal nodes marked either closed or open.
\end{definition}

\begin{itemize}
  \item Edges are local theorems or routing lemmas.
  \item Terminal nodes are either contradictions, lower already-closed strata, or explicitly open obstructions.
  \item A graph is useful even before every terminal node is closed, because it records exactly what remains to be proved.
\end{itemize}

\subsection{The abstract routed-closure principle}
\begin{theorem}[Abstract routed-closure principle]
\label{thm:abstract-routed-closure}
Let $G$ be a finite rooted routing graph.  Suppose every root branch enters at least one outgoing edge, every edge is justified by a valid proof capsule whose hypotheses are produced upstream, every retained obstruction is preserved along edges, and every terminal node is closed by contradiction with its retained obstruction.  Then no root branch exists.
\end{theorem}

\begin{proof}[Proof sketch]
Follow any hypothetical root branch through the finite graph; exhaustiveness gives a terminal node, preservation gives a surviving obstruction there, and terminal closure contradicts that obstruction.
\end{proof}

\subsection{Fault containment}
\begin{definition}[Fault containment]
A routed proof architecture has fault containment if invalidating one capsule changes only the closure status of the nodes and paths depending on that capsule, while leaving the rest of the graph meaningful.
\end{definition}

\begin{itemize}
  \item Every capsule has an area of effect: the topological region of the routing graph whose correctness genuinely depends on that local estimate.
  \item Fault containment is the demand that an error remain quarantined inside its area of effect rather than poisoning the proof trunk globally.
\end{itemize}

\begin{proposition}[Local invalidation principle]
If a capsule $e\in E(G)$ is invalidated, then all branches not passing through $e$ remain governed by the verified subgraph $G\setminus\{e\}$.
\end{proposition}

\begin{proof}[Proof sketch]
The proof is graph-theoretic: dependency is path-local, so invalidating an edge affects exactly those terminal claims whose proof paths use that edge.
\end{proof}

% =============================================================================
\section{Auditability and LLM-Oriented Proof Design}
\label{sec:auditability}

\subsection{Common failure modes of LLM-generated proofs}
\begin{itemize}
  \item Monolithic proofs exceed the active reasoning window of current language models and formal tactic engines, so relevant hypotheses fall out of scope while the search space still keeps growing.
  \item Once that context window is saturated, the tactic search space explodes: scaling conventions, compactness topologies, and gauge choices start interacting combinatorially instead of locally.
  \item The architecture addresses this by compressing each edge to a bounded local reasoning problem whose hypotheses, outputs, and failure routes fit inside one audit context.
  \item Hallucinated estimates still occur, but they now surface as local contract violations rather than as hidden defects dispersed across a hundred-page manuscript.
\end{itemize}

\subsection{LLM-proof generation and auditability}
\begin{itemize}
  \item The architecture acts as a compiler: it takes a Millennium-scale PDE problem and lowers it into a sequence of localized tactic puzzles that fit inside an LLM's active reasoning window.
  \item Each proof capsule is intentionally close to an undergraduate or first-year graduate exercise in local analysis, even when the surrounding theorem is globally deep.
  \item The architecture is therefore not only audit-friendly but search-friendly: it drastically shrinks the local tactic space that a model or proof assistant must explore.
  \item Axiomatizing the literature is part of the design.  Established theorems such as CKN concentration, Seregin-type limits, or downstream Liouville inputs can be imported as typed capsules while the machine verifies the routing between them.
  \item This turns AI assistance from vague theorem imitation into structured theorem transport across a graph of explicitly declared interfaces.
\end{itemize}

\subsection{Audit dimensions}
\begin{center}
\begin{longtable}{p{0.25\linewidth}p{0.65\linewidth}}
\toprule
Audit type & Local question \\
\midrule
Scaling audit & Do all terms transform with the claimed critical dimension? \\
Sign audit & Are modulation signs consistent across the chart, PDE, energy inequality, and cost identity? \\
Hypothesis audit & Is each theorem invoked only after its assumptions have been generated? \\
Locality audit & Does the estimate use only the declared cylinder, window, or compact state? \\
Pressure/gauge audit & Are pressure normalizations and gauge choices compatible under limits? \\
Compactness audit & Is the topology strong enough for every nonlinear passage? \\
Retention audit & Does the retained obstruction survive this edge? \\
Failure-route audit & If the estimate fails, is the branch routed to a named state? \\
Dependency audit & Is there any circular use of a closure theorem inside its own entry route? \\
\bottomrule
\end{longtable}
\end{center}

\subsection{Capsule prompt template for LLM audit}
\begin{lstlisting}
Task: Audit the following proof capsule.
Input contract:
  [List exact hypotheses.]
Output contract:
  [List exact conclusion.]
Retained obstruction:
  [Name the nonzero invariant.]
Failure routes:
  [List named exits if a hypothesis fails.]
Audit checks:
  1. Verify scaling of every term.
  2. Verify sign conventions.
  3. Verify pressure/gauge normalization.
  4. Verify compactness topology and nonlinear passage.
  5. Verify no hidden global estimate is used.
  6. Verify the retained obstruction survives.
  7. Verify all dependencies are upstream.
Return:
  PASS / FAIL / NEEDS ROUTING, with exact failing line or hypothesis.
\end{lstlisting}

\subsection{How fault containment helps human referees}
\begin{itemize}
  \item A referee can mark a capsule invalid without needing to reject the entire architecture.
  \item The graph then records a precise open node: missing estimate, unstable observable, nonexhaustive routing, or unclosed terminal stratum.
  \item This turns proof failure into useful diagnostic information rather than an opaque collapse.
  \item Allows to parallelize the repair process: different experts can work on different open nodes without stepping on each other's verified subgraph.
\end{itemize}

% =============================================================================
\section{The Navier--Stokes Program as a Case Study}
\label{sec:ns-case-study}

\subsection{The case-study role}
\begin{itemize}
  \item The \NS{} regularity problem is used as a stress test for the architecture because it combines scaling, pressure, compactness, and critical concentration.
  \item The Millennium-problem setting is ideal because it forces the architecture to confront the hardest interaction between nonlocal pressure, critical scaling, and compactness without hiding behind generic slogans.
  \item The case study should be read as a proof-design ledger whose analytic capsules are individually typed, auditable, and routable.
  \item Its value is not only explanatory: the routing graph is the logical AST that guarantees exhaustion and forbids ghost branches from escaping the local closure mechanisms.
\end{itemize}

\subsection{First gate: critical concentration}
\begin{itemize}
  \item Candidate singularity implies positive scale-critical concentration after excluding the epsilon-regularity branch.
  \item The retained obstruction begins as positive CKN or velocity concentration.
  \item Vanishing concentration is closed immediately by local epsilon-regularity.
\end{itemize}

\subsection{Second gate: Type I versus Type II}
\begin{itemize}
  \item Type I branch: local pointwise Type I envelope allows ancient-profile extraction.
  \item Type II branch: no local Type I envelope, so one studies retained local windows around moving concentration cores.
  \item The branches are separated because their compactness mechanisms are different.
\end{itemize}

\subsection{Type I architecture}
\begin{itemize}
  \item Pipeline: concentration $\to$ no-escape $\to$ Seregin extraction $\to$ bounded centered ancient profile $\to$ class routing.
  \item Direct closure rows: small, stationary $L^3$, uniformly tight, and structure/decay-covered ancient profiles.
  \item The point of the routing is that each closure uses only the local profile data generated by the branch, not a whole-space bound carried from the initial solution.
  \item Residual row: not a wastebasket, but a terminal-dynamics object with its own descendant, tail, affine, diffuse, and endpoint routes.
\end{itemize}

\subsection{Type II architecture}
\begin{itemize}
  \item Pipeline: retained core $\to$ repaired gauge $\to$ represented local equation $\to$ compactness/cost/multibubble/cascade/terminal routing.
  \item Repaired gauges turn scale drift and center drift into explicit modulation coefficients rather than hidden compactness errors.
  \item Type II routing discards global critical norms in favor of compact-window observables, local Calder\'on--Zygmund pressure control, and scale-collapse cost estimates tied to the active core.
  \item Cost and carrier routes treat scale collapse without assuming a global critical norm or a whole-space Liouville theorem.
\end{itemize}

\subsection{State-space diagram in words}
\begin{lstlisting}
terminal singular branch
  -> concentration gate
      -> no concentration: regularity
      -> positive concentration
          -> Type I envelope
              -> Seregin ancient state
              -> direct Liouville classes or residual class
          -> Type II branch
              -> repaired gauge / local state-space alternatives
              -> compact core, rough core, multibubble, cascade,
                 scale collapse, carrier, terminal profile, exterior exits
\end{lstlisting}

\begin{itemize}
  \item The flow diagram is not a mere summary figure; it is the explicit AST of the proof, recording exhaustion of cases and forbidding undeclared branch escapes.
  \item Once a branch enters this tree, every future estimate is attached to a named local state and every terminal trap is tied to a retained obstruction.
\end{itemize}

% =============================================================================
\section{Verification Ledgers}
\label{sec:ledgers}

\subsection{Interface ledger}
\begin{center}
\begin{tabular}{p{0.23\linewidth}p{0.31\linewidth}p{0.31\linewidth}}
\toprule
Interface & Source output & Target input \\
\midrule
Concentration gate & positive critical mass & Type I/Type II routing \\
Type I extraction & bounded ancient profile with retained mass & direct classes or residual class \\
Residual routing & realized descendant or terminal state & matching residual closure capsule \\
Type II entry & retained moving core & repaired-gauge representation \\
Gauge representation & local equation with $a,b$ and pressure atlas & compactness, Caccioppoli, cost, and carrier capsules \\
Terminal profile route & retained terminal mass & profile extraction or hidden-mass exit \\
\bottomrule
\end{tabular}
\end{center}

\subsection{Retained-obstruction ledger}
\begin{center}
\begin{longtable}{p{0.25\linewidth}p{0.35\linewidth}p{0.30\linewidth}}
\toprule
Obstruction & Produced by & Contradicted by \\
\midrule
positive CKN mass & concentration gate & epsilon-regularity or vanishing pressure/velocity \\
compact local $L^3$ mass & Seregin extraction & Liouville zero conclusion \\
residual ancestry & descendant realization & closure of non-realized auxiliary objects is disallowed \\
retained Type II core mass & selected core construction & compact zero limit or finite-cost collapse \\
carrier cost & local cost validation & corrected monotonicity plus carrier routing \\
terminal critical mass & terminal profile extraction & profile exhaustion and perturbative remainder \\
\bottomrule
\end{longtable}
\end{center}

\subsection{Closure-status ledger}
\begin{itemize}
  \item Every terminal node should have status \texttt{closed}, \texttt{imported}, \texttt{open}, or \texttt{conjectural}.
  \item Every imported theorem should be stated with exact hypotheses and bibliographic source.
  \item Every open node should be allowed to remain open without invalidating unrelated closed nodes.
\end{itemize}

\subsection{Repair protocol}
\begin{enumerate}
  \item Identify the failing capsule and its outgoing edge.
  \item Decide whether the estimate is false, under-hypothesized, or assigned to the wrong state.
  \item Add a named failure route if the negation has stable PDE meaning.
  \item Split the state if two geometrically different failures were conflated.
  \item Re-run the dependency audit for all downstream closures.
\end{enumerate}

% =============================================================================
\section{Portability Beyond Navier--Stokes}
\label{sec:portability}

\subsection{General principle}
\begin{itemize}
  \item The portable idea is not to copy the \NS{} strata.
  \item The portable idea is to identify invariant observables, route failure modes, and use only rigidity theorems whose hypotheses have been generated.
\end{itemize}

\subsection{Potential targets}
\begin{center}
\begin{longtable}{p{0.23\linewidth}p{0.36\linewidth}p{0.31\linewidth}}
\toprule
Problem class & Likely retained obstruction & Likely states \\
\midrule
Geometric flows & curvature concentration or entropy drop & tangent flows, necks, bubbles, ancient solitons \\
Dispersive PDE & critical mass/energy or scattering norm & profiles, soliton channels, radiation, cascades \\
Gauge theories & energy concentration and gauge class & bubbles, Coulomb gauges, neck regions, defect measures \\
Harmonic map heat flow & energy density and topological charge & bubble trees, necks, no-neck defects \\
Nonlinear wave equations & energy packets and channels & profiles, self-similar cones, exterior radiation \\
\bottomrule
\end{longtable}
\end{center}

\subsection{What changes in another PDE}
\begin{itemize}
  \item The critical observable changes.
  \item The compactness topology changes.
  \item The gauge or symmetry group changes.
  \item The closure theorems change.
  \item The architecture remains: retained obstruction, routing graph, proof capsules, terminal closures, fault containment.
\end{itemize}

% =============================================================================
\section{Exception Handling and Proof Resilience}
\label{sec:resilience}

\subsection{Dynamic repair via sub-branching}
\begin{itemize}
  \item When a reviewer or machine auditor finds an unclosed estimate, the architecture does not collapse; it instructs the proof to spawn a new named geometric stratum capturing the observed failure mode.
  \item A failed compactness route can become a compactness-defect branch, a failed modulation bound can become a modulation-failure branch, and a failed pressure normalization can become a pressure-atlas branch.
  \item The key principle is that unresolved behavior is promoted into a first-class state rather than left as an informal gap inside the trunk of the proof.
  \item Sub-branching is therefore not a sign of weakness but the mechanism by which the architecture converts critique into new local mathematics.
\end{itemize}

\subsection{Exception handling matrix}
\begin{center}
\begin{longtable}{p{0.32\linewidth}p{0.58\linewidth}}
\toprule
Exception detected & Resilience response \\
\midrule
scaling formula wrong & quarantine the chart capsule, invalidate only downstream sign-dependent edges, and preserve the rest of the routed graph \\
compactness lemma too strong & spawn a compactness-failure stratum or strengthen the input contract before re-entering the main branch \\
pressure gauge incompatible & repair the local pressure atlas or reroute to a pressure/exterior observer state \\
retained mass not preserved & replace the retained observable or split the branch until preservation is restored locally \\
Liouville theorem invoked too early & move the closure capsule downstream until its hypotheses are actually generated \\
cost identity wrong sign & reopen the scale-collapse branch and keep all unrelated branches certified \\
residual class not exhausted & mark the residual node open, attach new descendant routes, and leave upstream routing intact \\
\bottomrule
\end{longtable}
\end{center}

% =============================================================================
\section{Conclusion}
\label{sec:conclusion}

\begin{itemize}
  \item Long PDE regularity proofs should be organized as auditable routing graphs, not only as linear manuscripts.
  \item The decisive move is to discard hidden global estimates in favor of local state-space stratifications whose observables, gauges, and closure mechanisms remain bounded and machine-readable.
  \item Local proof capsules make LLM assistance stronger because every estimate is forced to declare hypotheses, outputs, retained obstructions, and failure routes inside one audit window.
  \item The monolithic human-readable proof is no longer the only viable endpoint for nonlinear mathematics; distributed, fault-tolerant, machine-auditable proof networks are becoming the natural format for problems of this complexity.
  \item By compiling continuous geometric mechanics into a discrete proof AST, the architecture bridges the gap between critical PDE analysis and automated theorem proving in systems such as Lean~4.
  \item The next step is to maintain a machine-readable ledger linking every capsule to its source theorem, dependency path, status, and audit history.
\end{itemize}

% =============================================================================
\appendix

\section{Proof-Capsule Template}
\label{app:capsule-template}

\subsection{Minimal template}
\begin{lstlisting}
Capsule name:
Implementation labels:
  L1. [theorem/proposition/lemma label]
  L2. [supporting theorem/proposition/lemma label]
Input contract:
  H1. [local cylinder/window/state]
  H2. [bounds and topology]
  H3. [pressure/gauge normalization]
Output contract:
  O1. [local estimate or routing conclusion]
Retained obstruction:
  I. [quantity and survival statement]
Failure routes:
  F1. [named exit]
  F2. [named exit]
Audit checklist:
  A1. scaling
  A2. sign convention
  A3. pressure/gauge
  A4. compactness topology
  A5. nonlinear passage
  A6. retention
  A7. dependencies
Closure status:
  verified / imported / open / conjectural
\end{lstlisting}

\subsection{Example capsule: selected-window compactness extraction}
\begin{itemize}
  \item Capsule name: Selected-window compactness extraction.
  \item Implementation labels:
  \item L1. \texttt{paper1:prop:compactness}.
  \item L2. \texttt{paper1:lem:D-localization} for the smaller-cylinder pressure bound used inside the compactness proof.
  \item Source: Proposition ``Compactness on selected windows'' in the companion manuscript \emph{Type II Regularity for Navier--Stokes}.
  \item Input contract:
  \item H1. A represented suitable sequence on $Q_{R_*}$ solving the local renormalized equation.
  \item H2. Uniform compact-cylinder bounds $A_{V_n}(R_*)+E_{V_n}(R_*)+C_{V_n}(R_*)+D_{P_n}(R_*)\le M$.
  \item H3. Uniformly bounded modulation coefficients $a_n,b_n$ on the selected time interval.
  \item Output contract:
  \item O1. For every $r<R_*$, a subsequence converges weak-* in $L^\infty_tL^2_x$, weakly in $L^2_tH^1_x$, and strongly in $L^3(Q_r)$.
  \item O2. Mean-normalized pressure converges weakly in $L^{3/2}(Q_r)$, and the limit still satisfies the represented equation distributionally on smaller cylinders.
  \item Retained obstruction:
  \item I. The capsule does not itself contradict the retained obstruction; it transports the branch into a compactness topology strong enough for later retention and contradiction capsules.
  \item Failure routes:
  \item F1. If the local pressure oscillation bound is unavailable, route to a pressure-row or pressure-compactness failure state.
  \item F2. If bounded modulation or the time-derivative control fails, route to modulation-failure or compactness-failure states rather than invoking Aubin--Lions implicitly.
  \item Audit checklist:
  \item A1. Check the pressure is normalized by spatial means on each compact ball.
  \item A2. Check the $L^2$ to $L^3$ interpolation and the use of the $L^{10/3}$ bound.
  \item A3. Check weak-* convergence for $a_n,b_n$ is sufficient for passage to the limit in the lower-order terms.
  \item A4. Check the argument never upgrades local compact-cylinder data into a hidden whole-space compactness theorem.
  \item Closure status:
  \item Open as a transport capsule; it is an upstream interface rather than a terminal contradiction.
\end{itemize}

\subsection{Example capsule: local single-core Type II exclusion}
\begin{itemize}
  \item Capsule name: Local single-core Type II exclusion.
  \item Implementation labels:
  \item L1. \texttt{paper1:thm:single-core-criterion}.
  \item L2. \texttt{paper1:lem:zero-dissipation} for the constant-limit reduction.
  \item L3. \texttt{paper1:lem:zero-contradiction} for the zero-limit contradiction.
  \item L4. \texttt{paper1:lem:nonzero-contradiction} for the bounded-window exit.
  \item Source: Theorem ``Local single-core Type II criterion'' in the companion manuscript \emph{Type II Regularity for Navier--Stokes}.
  \item Input contract:
  \item H1. A compact single-core vanishing-dissipation branch in the precise sense of the Type II manuscript.
  \item H2. Upstream access to the named inputs ``passage to selected-window branches,'' ``local representation input,'' and ``local pressure reconstruction.''
  \item H3. Persistence of a selected retained core with positive concentration on $Q_{r_*}$.
  \item Output contract:
  \item O1. The branch cannot remain a local Type II sequence.
  \item O2. Any nonzero bounded selected-window limit is rerouted to the bounded-window row instead of remaining inside the strict selected-window branch.
  \item Retained obstruction:
  \item I. Positive retained velocity concentration on the selected core, encoded by a fixed lower bound $C_{V_n}(r_*)\ge \eta_0$ along the selected-window subsequence.
  \item Failure routes:
  \item F1. If the limit is spatially constant and nonzero, exit the strict selected-window Type II branch through the bounded-window routing row.
  \item F2. Multibubble, cascade, gauge-degenerate, and compactness-failure alternatives are not hidden inside the theorem; they are separate named rows outside the single-core capsule.
  \item Audit checklist:
  \item A1. Check the vanishing localized dissipation really implies a spatially constant limit on the selected compact cylinder.
  \item A2. Check the strong $L^3$ convergence on the core cylinder is exactly what forces the zero-limit contradiction.
  \item A3. Check the retained obstruction is contradicted only in the zero-constant case and rerouted rather than contradicted in the nonzero-constant case.
  \item A4. Check no hidden stationary, ancient, or self-similar Liouville theorem is invoked.
  \item Closure status:
  \item Verified terminal exclusion capsule inside the local single-core row, with named exits for behavior that leaves that row.
\end{itemize}

\section{LLM Audit Prompt Library}
\label{app:prompt-library}

\subsection{Scaling audit prompt}
\begin{lstlisting}
Given the change of variables and PDE below, verify every scaling exponent.
Return a table with columns: term, original scaling, transformed scaling, expected coefficient.
Flag missing powers or inconsistent conventions.
\end{lstlisting}

\subsection{Sign audit prompt}
\begin{lstlisting}
Audit the sign convention for the modulation coefficients a and b.
Check consistency among: chart definition, chain rule, represented PDE,
local energy inequality, Caccioppoli estimate, and cost identity.
Return the first inconsistent formula if any.
\end{lstlisting}

\subsection{Hypothesis audit prompt}
\begin{lstlisting}
For each theorem invocation in this proof segment, list the invoked theorem,
its hypotheses, the upstream lines that produce each hypothesis, and any missing hypothesis.
Classify each invocation as VALID, MISSING HYPOTHESIS, or CIRCULAR.
\end{lstlisting}

\subsection{Pressure/gauge audit prompt}
\begin{lstlisting}
Check whether pressure is compared only modulo functions of time or spatial means.
Verify that every pressure compactness or pressure-flux step uses the declared normalization.
Flag any absolute pressure estimate not invariant under pressure gauge changes.
\end{lstlisting}

\subsection{Retention audit prompt}
\begin{lstlisting}
Identify the retained obstruction entering this capsule.
Track it through subsequences, rescalings, gauge changes, and limits.
State exactly where the terminal contradiction uses the retained obstruction.
\end{lstlisting}

\section{Navier--Stokes Routing Graph Skeleton}
\label{app:ns-routing-graph}

\subsection{Global routing}
\begin{lstlisting}
finite-energy suitable weak solution
  -> suppose terminal singular point exists
      -> concentration gate
          -> no concentration: epsilon-regularity closure
          -> positive concentration
              -> Type I envelope?
                  yes: Seregin ancient-profile state space
                  no: Type II retained local-window state space
\end{lstlisting}

\subsection{Type I routing}
\begin{lstlisting}
Type I branch
  -> Seregin extraction
      -> small class: perturbative Liouville
      -> stationary L3 class: stationary rigidity
      -> uniformly tight class: endpoint ancient Liouville
      -> structure/decay class: specialized Liouville theorem
      -> residual class: residual terminal-dynamics graph
\end{lstlisting}

\subsection{Residual Type I routing}
\begin{lstlisting}
residual class
  -> axisymmetric / rotational / stationary-hull / affine-parasitic strata
  -> tail strata: log-diffuse, Young defect, coherent critical tails
  -> terminal strata: nonactive, diffuse, mixed, active descendants
  -> generic residual: parasitic-free mildness + sequence L3 + endpoint Liouville
\end{lstlisting}

\subsection{Type II routing}
\begin{lstlisting}
Type II branch
  -> select retained moving core
      -> repaired gauge representation
      -> compact single core
      -> rough core / H1 loss
      -> multibubble / cascade / gauge degeneracy
      -> scale collapse / finite cost / carrier routing
      -> terminal critical profile / exterior states
\end{lstlisting}

\section{Verification Ledger Templates}
\label{app:ledger-templates}

\subsection{Closure status table template}
\begin{center}
\begin{longtable}{p{0.18\linewidth}p{0.22\linewidth}p{0.17\linewidth}p{0.25\linewidth}}
\toprule
Node & Closure mechanism & Status & Audit notes \\
\midrule
{[Node]} & {[Capsule]} & verified/open & {[Notes]} \\
\bottomrule
\end{longtable}
\end{center}

\subsection{Dependency ledger template}
\begin{center}
\begin{longtable}{p{0.20\linewidth}p{0.22\linewidth}p{0.22\linewidth}p{0.20\linewidth}}
\toprule
Capsule & Uses & Produces & Downstream consumers \\
\midrule
{[Capsule]} & {[Inputs]} & {[Outputs]} & {[Downstream]} \\
\bottomrule
\end{longtable}
\end{center}

\subsection{Retained-obstruction ledger template}
\begin{center}
\begin{longtable}{p{0.22\linewidth}p{0.22\linewidth}p{0.22\linewidth}p{0.16\linewidth}}
\toprule
Obstruction & Produced at & Preserved through & Terminal use \\
\midrule
{[Invariant]} & {[Produced at]} & {[Edges]} & {[Use]} \\
\bottomrule
\end{longtable}
\end{center}

\section{Expansion Plan for the Full Manuscript}
\label{app:expansion-plan}

\subsection{Main text expansion targets}
\begin{itemize}
  \item Expand \Cref{sec:introduction} into 4--6 pages with motivating examples of fragile monolithic proofs.
  \item Expand \Cref{sec:global-bottleneck} into the conceptual bridge between local proof design and the elimination of whole-space critical estimates.
  \item Expand \Cref{sec:abstract-architecture} into the formal mathematical core of the paper.
  \item Keep \Cref{sec:auditability} short, practical, and free of hype.
  \item Keep \Cref{sec:ns-case-study} as a case study with explicit disclaimers about analytic verification.
  \item Put all large theorem-register material in appendices or supplementary files.
\end{itemize}

\subsection{Professional style rules}
\begin{itemize}
  \item Avoid saying ``LLMs prove the theorem''; say ``LLMs audit local capsules.''
  \item Avoid saying ``all exits are closed'' unless every closure capsule has been independently checked.
  \item Use ``open node'' instead of ``gap'' when the architecture intentionally records an unclosed branch.
  \item Use ``proof capsule'' only for statements with explicit hypotheses and outputs.
  \item Use ``case study'' for \NS{} until the analytic proof series is independently mature.
\end{itemize}

\subsection{Suggested companion artifacts}
\begin{itemize}
  \item A machine-readable JSON/YAML ledger of capsules and dependencies.
  \item A theorem-label register generated from the TeX source.
  \item A prompt library for local audits.
  \item A status dashboard marking capsules as verified, imported, open, or under revision.
\end{itemize}

\begin{thebibliography}{99}

\bibitem{CKN82}
L. Caffarelli, R. Kohn, and L. Nirenberg.
Partial regularity of suitable weak solutions of the Navier--Stokes equations.
\emph{Communications on Pure and Applied Mathematics}, 35(6):771--831, 1982.

\bibitem{Leray34}
J. Leray.
Sur le mouvement d'un liquide visqueux emplissant l'espace.
\emph{Acta Mathematica}, 63:193--248, 1934.

\bibitem{ESS03}
L. Escauriaza, G. Seregin, and V. Sverak.
$L_{3,\infty}$-solutions of Navier--Stokes equations and backward uniqueness.
\emph{Russian Mathematical Surveys}, 58(2):211--250, 2003.

\bibitem{NRS96}
J. Necas, M. Ruzicka, and V. Sverak.
On Leray's self-similar solutions of the Navier--Stokes equations.
\emph{Acta Mathematica}, 176:283--294, 1996.

\bibitem{AubinSimon}
J.-P. Aubin and J. Simon.
Compactness tools for evolution equations.
\emph{[Replace with exact references in final bibliography]}.

\bibitem{FormalMethodsPlaceholder}
[Placeholder]
References on formal proof assistants, proof engineering, or LLM-assisted theorem proving should be added here only after exact bibliographic verification.

\end{thebibliography}

\end{document}
